% 2.1.1
% Ficheiro incluído por main.tex

\subsection[Preparar o cliente: confirmar que o Winbind está desativado]{Preparar o cliente: confirmar que o Winbind está desativado}

\passo{Passo 1 - Confirmar que o serviço Winbind está desativado}

\begin{lstlisting}[style=bashstyle]
systemctl is-enabled winbind
\end{lstlisting}

\emph{Verificação esperada:} \texttt{disabled}. Se aparecer \texttt{enabled}: \texttt{sudo systemctl disable --now winbind}.

\passo{Passo 2 - Confirmar que o Winbind não está no PAM}

\begin{lstlisting}[style=bashstyle]
grep -c pam_winbind /etc/pam.d/common-auth
\end{lstlisting}

\emph{Verificação esperada:} \texttt{0}. Se der outro número: \texttt{sudo pam-auth-update --disable winbind}.

\passo{Passo 3 - Confirmar que o Winbind não está no NSS}

\begin{lstlisting}[style=bashstyle]
grep -E '^(passwd|group):' /etc/nsswitch.conf
\end{lstlisting}

\emph{Verificação esperada:} as linhas \texttt{passwd:} e \texttt{group:} não têm \texttt{winbind} (ex.: \texttt{passwd: files systemd sss}). Se tiverem: \texttt{sudo sed -i -E '/\allowbreak{}\textasciicircum{}(passwd\textbar{}group):/\allowbreak{} s/\allowbreak{} winbind/\allowbreak{}/\allowbreak{}' /\allowbreak{}etc/\allowbreak{}nsswitch.conf}.

